\documentclass[10pt,a4paper]{article}

\usepackage[T1]{fontenc}
\usepackage[utf8]{inputenc}
\usepackage{lmodern}
\usepackage[margin=20mm,top=18mm,bottom=20mm]{geometry}
\usepackage{booktabs}
\usepackage{enumitem}
\usepackage{titlesec}
\usepackage{xcolor}
\usepackage[hidelinks]{hyperref}

\definecolor{rule}{gray}{0.55}

\titleformat{\section}{\large\bfseries}{\thesection.}{0.6em}{}
\titlespacing*{\section}{0pt}{1.4ex plus .3ex}{0.7ex}
\titleformat{\subsection}{\normalsize\bfseries}{}{0em}{}
\titlespacing*{\subsection}{0pt}{1.1ex plus .2ex}{0.4ex}

\setlength{\parindent}{0pt}
\setlength{\parskip}{0.55ex plus 0.2ex}
\setlist[itemize]{leftmargin=1.1em,itemsep=0.25ex,topsep=0.35ex,parsep=0pt}
\setlist[description]{leftmargin=0pt,itemsep=0.5ex,topsep=0.5ex,parsep=0pt,style=unboxed,font=\bfseries}

\newcommand{\term}[1]{\textit{#1}}

\title{\vspace{-8mm}\textbf{\large In-Memory Account Ledger Core}\\[2pt]
\normalsize Architecture and Trade-offs}
\author{\normalsize Dmitry Bannikov}
\date{\normalsize 10 September 2026}

\begin{document}
\maketitle
\vspace{-6mm}

\section{Append-only at scale}

\subsection{What breaks first at 100$\times$ volume}

Not memory. A six-day window of entries is small even at a hundred times the
volume, and the log is the cheapest thing in the system to store. What breaks
is the \term{daily close}, and it breaks quadratically.

Every balance in this design is derived by scanning the log for entries with
$\textit{value\_date} \le d$. That is $O(N)$ in the number of entries. The
close calls it once per account per candidate day, and under the retroactive
fee policy the candidate set is every day back to the start of the window, so
the close is $O(D \cdot A \cdot N)$. Since $N$ itself grows with volume,
multiplying throughput by 100 multiplies close-time work by roughly $10^4$. A
close that takes seconds today takes hours, and it takes them inside the
overnight window where there is no slack.

The second thing to break is worse, because it is on a latency budget rather
than a batch window: the available-balance check on the authorisation path.
Deriving a hold's current state means scanning the hold log for that
authorisation, and summing active holds means doing that once per
authorisation the account has ever had. The cost of approving a card
transaction therefore grows with the account's \term{lifetime} authorisation
count, not with the number of holds actually outstanding. An account that has
been open a year is slower to authorise than one opened yesterday holding the
same funds. That is the wrong quantity to be sensitive to, and it degrades
silently until it starts timing out at the scheme.

\subsection{Where unbounded state accumulates}

The log grows without bound, and that is correct: it is the record, and I would
not bound it. The defect is not unbounded storage, it is \term{unbounded read
amplification} --- three places where the cost of answering a bounded question
grows with total history:

\begin{itemize}
  \item Every balance query reads from the beginning of time. There is no lower
        bound on the scan, so answering ``what is the balance today'' gets more
        expensive every day, forever.
  \item Settled, released and declined authorisations never leave the hot path.
        They are dead records that every subsequent approval decision still
        pays to skip.
  \item The close-time snapshot map grows one entry per day per account, and
        nothing ever reads the old ones except a report.
\end{itemize}

\subsection{The cheapest change that defers it}

In two tiers, because the cheap fix and the right fix are not the same fix.

\textbf{Immediately, and for almost nothing:} maintain a per-(account, value
date) delta bucket on append and derive balances by prefix-summing at most $D$
buckets rather than $N$ entries. A backdated entry invalidates only the prefix
sums at or after its value date. This turns the close from quadratic in volume
into linear in days and touches no semantics, because the index is still
\term{derived} --- it can be rebuilt from the log, and a test that asserts the
index equals a full replay is a property the design already needs for other
reasons. Partitioning holds into open and closed sets is the same kind of
change and removes the authorisation-path problem entirely.

\textbf{Structurally:} materialise a cumulative balance checkpoint per account
at period boundaries and read \term{checkpoint plus delta}. The obstacle is
that a checkpoint is only sound if the period behind it cannot change, and
value-dating is precisely the mechanism by which it can. So the checkpoint is
not really an indexing change --- it is a policy change wearing an indexing
costume, and the policy is a \term{correction window}: past some cut-off, an
entry value-dated into a closed period is refused at intake and rebooked as an
adjustment in the open period, carrying the original date as data rather than
as position.

That is the change I would argue for first, because it is the only one that
pays twice. It bounds the scan, it makes the ledger reconcilable against a
general ledger whose periods are already hard-closed, and it bounds the
retroactive recomputation that my own failing test exposes --- where the fees a
customer pays depend on the order corrections arrived in rather than on the
contents of the log. One control, two classes of problem.

\section{Value-dated entries in production}

\subsection{Operational surface}

A value-dated entry is a licence to rewrite a day that has already been
reported, reconciled, and in most cases sent to the customer. Everything
downstream of a closed day becomes mutable with it.

\begin{itemize}
  \item \textbf{Statement restatement.} A statement is a customer-facing
        document with legal weight. A backdated entry makes an issued statement
        wrong, which means the bank needs statement versioning, correction
        advices, and the ability to reproduce a statement exactly as it was
        issued --- not merely as the ledger reads now. This is why the report
        in Part~1 keeps both the close-time figure and the restated one; the
        pair is the minimum needed to explain a charge to the person who
        received it.
  \item \textbf{Prior-period adjustment.} The general ledger closes on a
        calendar; the core ledger, as built, does not. A backdated entry that
        crosses a month end lands in a GL period that no longer accepts
        postings, and one that crosses a year end touches audited financials.
        Without an explicit adjustment path the two ledgers simply diverge, and
        the difference surfaces as an unexplained suspense balance.
  \item \textbf{Retroactive charges.} The highest-risk output of the whole
        mechanism. The CBUAE Consumer Protection Regulation (Circular 8/2020)
        and its accompanying Standards require charges to be disclosed,
        transparent and explainable, with a complaints route. A fee levied for
        a day whose issued statement showed the account in credit cannot be
        explained from any document the customer holds.
  \item \textbf{Tax periods.} Explicit fee income carries 5\% VAT. A fee dated
        into an already-filed period changes a filed return.
  \item \textbf{AML and monitoring.} Transaction monitoring aggregates by date;
        a backdated entry can retrospectively change a threshold aggregation
        that a structuring rule already evaluated. Under Federal Decree-Law
        20/2018 the record must also survive five-year retention in both its
        original and corrected form --- the correction does not replace the
        original, and the system has to be able to produce both.
  \item \textbf{Internal fraud and earnings management.} Whoever chooses a
        value date chooses the fee, the interest and the reporting period into
        which revenue falls. Backdating is a control surface before it is a
        feature, and in the model as built no event records who chose.
\end{itemize}

\subsection{One control before go-live}

An \textbf{enforced value-date correction window with four-eyes approval and
mandatory reason coding, applied at the ledger boundary.}

An entry whose value date is older than a published cut-off is rejected at
intake. Booking one anyway requires a second authorised approver, a reason code
drawn from a closed list, and it automatically raises a customer notification
whenever the entry changes an issued statement or levies a charge.

I would choose this over the obvious alternatives --- a nightly restatement
reconciliation report, an alerting threshold on backdated volume, a shadow GL
--- for three reasons. It is preventive rather than detective, so the exposure
never opens rather than being discovered afterwards. It puts a named human
against every retroactive charge, which is the first thing both a
consumer-protection complaint and a supervisory examination ask for. And it is
the same mechanism that makes checkpointing sound in \S1, so it is one control
answering a regulatory question and a scaling question at once.

\section{Authorisation lifecycle}

A matching settlement is the least common way an authorisation ends. Every
other ending below is reachable in this domain, and each needs a mandated
behaviour rather than a default.

One rule governs all of them, and the model as built already follows it: an
ending is always an \term{event}, never a background sweep that mutates state.
If available balance rises and no event says why, nobody can answer the
customer's question. A second rule follows from the first: what is released is
always the \term{remaining} hold, computed at that moment --- never the
original amount added back.

\begin{description}
  \item[Declined at authorisation.] Insufficient available balance, a risk
        decision, or an invalid card state. \emph{Behaviour:} record the
        decline with its reason and the available balance at that instant, and
        keep it queryable. Declines are a top complaint category and a fraud
        signal; a decline that leaves no trace is an unanswerable support
        ticket. No hold is created.
  \item[Expiry without presentment.] The acquirer never presents, or presents
        after the hold window --- an abandoned checkout, or a pre-authorisation
        the merchant simply never uses. \emph{Behaviour:} release automatically
        at a scheme- and merchant-category-specific window, as an explicit
        expiry event. A settlement arriving afterwards is not an error; it
        becomes a force post and is handled as one.
  \item[Reversal by the acquirer, full or partial.] The terminal cancels, or an
        estimated authorisation is reduced to the amount actually transacted
        --- the fuel pump that authorises an estimate and reverses the unused
        portion. \emph{Behaviour:} release exactly the reversed amount and leave
        the authorisation live at the reduced figure. Treating a partial
        reversal as a full one over-releases funds; treating it as invalid
        strands them.
  \item[Partial settlement, residual never presented.] A split shipment where
        the remainder is cancelled. \emph{Behaviour:} each settlement posts and
        decrements the hold; the authorisation ends when the hold reaches zero
        or the expiry window closes, whichever comes first --- not on the first
        settlement.
  \item[Account frozen, closed, or the customer deceased.] A sanctions match, a
        court order, or an estate. \emph{Behaviour:} the hold must not simply be
        released, because releasing it makes the funds available for
        withdrawal. Move the held amount to a blocked position and let the
        eventual settlement or expiry resolve against that, not against the
        customer's available balance.
  \item[Stand-in or offline authorisation the issuer never saw.] The scheme
        approved on the issuer's behalf during an outage, or the terminal
        authorised below its floor limit. \emph{Behaviour:} there is no
        authorisation to end --- the ledger learns of the obligation only at
        clearing, and must accept it as a force post rather than reject a
        payment the scheme has already guaranteed. Matching against the
        clearing reference, with an amount tolerance, is what separates this
        from a genuinely bad message.
\end{description}

\section{What I cut, and what each cut defers}

\begin{description}
  \item[No intake idempotency.] The same event replayed posts twice. Every
        realistic upstream --- a scheme feed, a queue, a retrying client ---
        delivers at least once, so this is a live double-charge defect rather
        than a theoretical one. It is the first thing I would fix, and it is
        cheap: a unique constraint on the event identifier, enforced by the
        store rather than by a read-then-write check.
  \item[No concurrency control.] Single-threaded, with a sequence counter that
        is not atomic. Two writers interleaving would make log order stop
        matching audit order, which quietly destroys the one property the whole
        design rests on. Production needs a single writer per account or an
        explicit ordering service --- not a lock added later.
  \item[Fees derived rather than event-sourced.] Fee assessment recomputes from
        current balances instead of being an event carrying the as-of watermark
        it was computed against. This is the flaw the failing test documents:
        identical entries yield different fees depending on arrival order, so
        any rebuild or replica that re-derives fees disagrees with production.
        Deferred risk is a reconciliation break that appears only in disaster
        recovery, which is the worst time to find it.
  \item[No actor identity on any event.] Nothing records who booked an entry or
        chose its value date. This defers exactly the control \S2 depends on:
        four-eyes approval is unimplementable on a model that has no notion of
        eyes.
  \item[One settlement per authorisation; no expiry; no partial reversal.]
        Three of the six endings in \S3 are unrepresentable. Deferred risk is
        customer funds held indefinitely against authorisations that will never
        settle, and split-shipment settlements rejected as duplicates.
  \item[No fee-reversal policy.] When the entry that caused a fee is reversed,
        the fee stands, because no policy was given and inventing one silently
        seemed worse than leaving it visible. Deferred risk is a
        consumer-protection exposure attached to an undocumented decision ---
        the kind that is discovered by a complaint rather than by a review.
  \item[Single-currency fee schedule, no FX.] A fee owed on an account whose
        currency the schedule does not name raises rather than guesses at a
        rate. Deferred risk is a hard failure on a path that a second currency
        pair would reach immediately; the fix is a per-currency schedule, not a
        conversion.
  \item[$O(N)$ balance derivation and no checkpointing.] Deliberate, to keep
        the log the single source of truth while the design was still moving.
        Deferred risk is \S1 in full.
\end{description}

\end{document}
